\documentclass[11pt,a4paper]{article}
\usepackage[margin=2.5cm]{geometry}
\usepackage{amsmath,amssymb,booktabs,siunitx,graphicx,microtype}
\usepackage[colorlinks=true,allcolors=blue]{hyperref}
\usepackage{xcolor}
\usepackage[version=4]{mhchem}
\usepackage{abstract}
\usepackage{titlesec}
\usepackage{float}
\usepackage{placeins}
\usepackage[font=small,labelfont=bf,skip=6pt]{caption}

\graphicspath{{figures/}{./}}

\titleformat{\section}{\normalfont\large\bfseries}{\thesection}{1em}{}
\titleformat{\subsection}{\normalfont\normalsize\bfseries}{\thesubsection}{1em}{}

\sisetup{detect-weight=true, detect-family=true}

\title{\textbf{Orbital Layout in Heavy-Hex LUCJ Sampling for\\
Sample-Based Quantum Diagonalization}}
\author{Maxim Solovev\\
\small RIKEN Center for Computational Science\\
\small Supervisors: Prof.\ S.\ Yunoki, Dr.\ T.\ Shirakawa}
\date{\today}

\begin{document}
\maketitle

\begin{abstract}
\noindent
In sample-based quantum diagonalization (SQD), the local unitary cluster
Jastrow (LUCJ) circuit acts as a sampler whose Jastrow couplings are
truncated by a heavy-hex locality mask. Which couplings survive is determined
entirely by the orbital-to-qubit assignment --- a quantity inherited from the
active-space construction without physical justification, and treated as a
fixed convention throughout the literature since the ansatz was introduced.
We show that this assignment changes the SQD subspace error by up to a factor
of eight at fixed mask, shot and determinant budgets, on two systems, with an
effect far exceeding sampling noise (between/within variance ratios of 27.2
and 278.2).

The controlling mechanism is \emph{subspace capture} --- the fraction of the
exact wavefunction contained in the sampled product subspace. Its correlation
with subspace error lies between $-0.88$ and $-0.98$ across twenty-one
distinct same-spin chains, three molecules, two orbital bases, two
correlation regimes and two search axes, including twelve chains held out
from all development, with no failures.

The central structural result is that the mask exposes two \emph{separable
but interacting} degrees of freedom: the same-spin \emph{chain ordering}
($N_{\text{orb}}!$ choices) and the opposite-spin \emph{anchor selection}
($\binom{N_{\text{orb}}}{N_{\text{orb}}/4}$ choices). At $N_{\text{orb}}=10$
these span comparable observed error ranges --- \SI{286}{\milli\hartree} over
50 random permutations and \SI{234}{\milli\hartree} over all 120 anchor
triples at a fixed chain --- from search spaces differing by four orders of
magnitude. Optimising anchors at fixed ordering compresses the spread across
twenty chains and reverses their ranking.

We prove that no purely opposite-spin quantity can select the optimal anchor
set at every chain: such quantities are exactly invariant under the same-spin
permutation, while the optimum is not. Three families of chain-aware
alternatives, pre-registered across three separate attempts, all fail to
improve on the simple opposite-spin retained Jastrow weight $S_0$. Yet under
the objective that matters in practice --- shortlisting rather than blind
selection --- $S_0$ performs extremely well: ranking all 120 candidates and
evaluating only the top five reaches a median and 90th-percentile normalised
regret of $0.002$ across the twelve held-out chains, against $0.181$ and
$0.621$ for random selection. A circuit-resource audit against a heavy-hex
coupling map confirms that better anchor sets require \emph{fewer} two-qubit
gates, not more.
\end{abstract}

\tableofcontents
\newpage

\section{Problem statement}

The heavy-hex locality mask retains same-spin Jastrow elements only at
nearest-neighbour positions $(p,p+1)$ together with the on-site diagonal
$(p,p)$, and opposite-spin elements only at $p \equiv 0 \bmod 4$. At
$N_{\text{orb}}=10$ this keeps 19 of 55 same-spin and 3 of 55 opposite-spin
couplings; at $N_{\text{orb}}=36$ it keeps 5.6\% of same-spin pairs.

The mask itself is dictated by hardware. What is \emph{not} dictated is the
assignment of logical molecular orbitals to the fixed hardware-compatible
positions, and that assignment determines which couplings survive. It is
inherited from the active-space construction without physical criterion.

To our knowledge no published work varies it or reports its effect on
accuracy. The zig-zag pattern introduced with the ansatz has been reused
unchanged throughout the subsequent literature, and where circuit layout is
discussed the concern is hardware quality --- avoiding high-error physical
qubits --- rather than which orbital occupies which position. The nearest
related work addresses the same starting problem from the opposite direction,
re-optimising the ansatz parameters to compensate for a fixed mask rather
than choosing what the mask retains.

Permutations act on the operator,
\begin{equation}
e^{\hat K} e^{i\hat J} e^{-\hat K}
= \bigl(e^{\hat K}P^{\top}\bigr)\, e^{i(PJP^{\top})}\, \bigl(Pe^{-\hat K}\bigr),
\end{equation}
leaving the Hamiltonian and the reference determinant untouched; the
untruncated ansatz is bit-identically invariant under $P$. It is masking
\emph{after} permutation that makes orderings differ.

Prior treatment --- including our own earlier work --- has regarded the
orbital-to-qubit assignment as a single object. It is not. A permutation
determines two logically separable things: which orbital pairs are adjacent
on the chain, fixing the same-spin mask, and which orbitals land at positions
$p \equiv 0 \bmod 4$, fixing the opposite-spin anchors. Because the chain mask
is invariant under reversal, these can be varied independently, although
Section~\ref{sec:compression} shows their effects interact and partially
compensate. Separating them is the organising principle of this report.

A further distinction matters throughout. In VQE the ansatz state \emph{is}
the answer, so minimising $\langle\psi\lvert\hat H\rvert\psi\rangle$ is the
objective. In SQD the ansatz is only a sampler: its bitstrings select a
determinant subspace which is then diagonalised exactly. A layout with a
slightly worse variational energy may therefore produce a better finite
subspace at fixed budget. Sections~\ref{sec:ansatz}
and~\ref{sec:transmission} show this is not hypothetical.

\section{Methods}

\subsection{Systems}

\begin{table}[htbp]\centering\small
\begin{tabular}{llll}
\toprule
& \ce{N2} & \ce{H10} chain & \ce{Cr2}\\
\midrule
Basis & 6-31G & STO-6G & cc-pVDZ\\
Active space & CAS(6,10) & CAS(10,10) & CAS(12,12)\\
Geometry & $R = \SI{1.55}{\angstrom}$ & $R = \SI{1.6}{\angstrom}$ & $R = \SI{1.68}{\angstrom}$\\
Orbitals & canonical & block-Boys localised & block-Boys localised\\
Qubits & 20 & 20 & 24\\
CI dimension & 14{,}400 & 63{,}504 & 853{,}776\\
Determinants per sector & 15 & 15 & 55\\
Subspace fraction of CI & 1.56\% & 0.354\% & 0.354\%\\
Ideal capture at budget & 0.9866 & 0.7554 & 0.881\\
Shots & $10^{6}$ & $2\times10^{6}$ & $2\times10^{6}$\\
\bottomrule
\end{tabular}
\caption{The three benchmark systems, all exactly diagonalisable. The
determinant budget on \ce{Cr2} was chosen to match \ce{H10}'s subspace
fraction of the CI space exactly, so that the systems are compared at
equivalent relative budgets rather than equivalent absolute ones.}
\end{table}

$R=\SI{1.8}{\angstrom}$ was the original \ce{H10} design point and proved
unusable: CCSD genuinely diverges there, with $E_{\text{corr}}$ oscillating
between \SIrange{-0.4}{-0.9}{\hartree} over 300 cycles. The benchmark was
moved to \SI{1.6}{\angstrom}, where CCSD converges in 48 cycles. The move was
forced by a reference-method breakdown, not by the results at either
geometry. On \ce{Cr2}, by contrast, active-space CCSD converged cleanly in 26
cycles ($E_{\text{corr}} = \SI{-0.610}{\hartree}$); we note this as an
observation about the active-space problem, not as a claim about \ce{Cr2}'s
multireference character in the full space.

\subsection{Pipeline}

PySCF produces the RHF reference, CCSD amplitudes, active-space FCIDUMP and
exact CASCI vector. Amplitudes initialise an \texttt{ffsim} LUCJ operator; the
locality mask is applied; Qiskit Aer samples bitstrings; the top determinants
per spin sector by marginal weight form a product subspace; and \texttt{sbd}
(Shirakawa's C++ implementation, used unmodified) performs the selected-basis
diagonalization.

Two parameter choices were established empirically and are load-bearing.
Full-rank double factorisation is required: at low rank the ansatz collapses
back toward Hartree--Fock, giving \SI{160}{\milli\hartree} error against CCSD
on \ce{H10} where full rank gives \SI{1.2}{\milli\hartree}. The
compressed-factorisation optimiser is disabled, since it was observed to
worsen energies and to introduce nondeterminism that would confound ordering
comparisons.

Block-Boys localisation is applied separately to the occupied-active and
virtual-active blocks. Localising over the full active space would mix
occupied with virtual orbitals, destroying the HF reference and leaving the
CCSD amplitudes undefined as occupied-to-virtual objects; no valid amplitude
transformation exists in that case. Under block localisation the
transformation is exact,
\begin{equation}
t^{2\prime}_{IJAB} = \sum_{ijab} U^{o}_{iI}U^{o}_{jJ}U^{v}_{aA}U^{v}_{bB}\,
t^{2}_{ijab},
\end{equation}
because the CCSD wavefunction $e^{T}\lvert\mathrm{HF}\rangle$ is invariant
under occupied-occupied and virtual-virtual rotations when the amplitudes are
rotated with them. CCSD is not re-run in the localised basis: PySCF's CCSD
assumes a diagonal Fock matrix for the amplitude denominators, which
localised orbitals do not provide.

\subsection{Reproducibility protocol}

Identical code executed in separate processes was found to produce different
FCIDUMP files and energies differing by \SIrange{8}{17}{\milli\hartree}, from
thread-dependent floating-point variation in RHF/CCSD. Because sampled
marginals decay steeply, determinant selection at tight budgets is
near-degenerate and numerical noise flips which strings are retained. All
results here therefore use reference data (RHF, CCSD, CASCI, FCIDUMP)
computed once and cached, with OMP/MKL/OpenBLAS thread counts pinned to one,
a fresh seeded simulator per evaluation, a single FCIDUMP shared across every
layout within an experiment, and incremental persistence of results.

Sampling stability is budget-dependent in a way worth stating precisely. The
analytic marginal ratio $w_{16}/w_{15}$ at the selection boundary is 0.504
for \ce{N2} but 0.989 for \ce{H10}. An early \ce{H10} run at
$5\times10^{5}$ shots showed \SI{43}{\milli\hartree} of seed scatter; at
$2\times10^{6}$ shots all five seeds return bit-identical energies and
determinant sets. A ratio near unity means ``requires more shots'', not
``unresolvable'': the absolute probability gap is fixed while discrimination
signal-to-noise grows as $\sqrt{N_{\text{shots}}}$.

\FloatBarrier
\section{Validation}

\begin{itemize}\itemsep3pt
  \item \textbf{Group benchmark reproduced.} The \ce{N2} benchmark bundled
        with \texttt{sbd} (1387 determinants at the $10^{-4}$ threshold)
        gives $-109.0483526946501$ Ha against the published
        $-109.04835269$ Ha, agreeing to twelve significant figures. This was
        re-run fresh for this report rather than quoted from an earlier log.
  \item \textbf{Limiting cases.} \ce{H2}/STO-3G returns RHF
        ($-1.116759307396$ Ha) from an HF-only determinant file and FCI
        ($-1.137283834489$ Ha) from the full sampled set.
  \item \textbf{Unmasked-permutation invariance.} With the mask disabled a
        permutation is a pure relabelling, so the SQD energy must be exactly
        invariant. Four orderings on \ce{N2} return bit-identical
        $-108.823644563978$ Ha with unit Jaccard overlap of the selected
        determinant sets. This validates the operator-permutation,
        bitstring-to-determinant and FCIDUMP-labelling chain end to end.
  \item \textbf{Two independent implementations agree exactly.} Mechanism~A
        permutes the operator ($PJP^{\top}$ absorbed into the orbital
        rotations) then applies a fixed mask; mechanism~B leaves the operator
        untouched and selects which physical orbital pairs enter the
        interaction-pair list. Entrywise comparison of the resulting
        diag-Coulomb matrices and orbital rotations gives exact
        floating-point zero across all layouts tested.
  \item \textbf{The truncation is a projection, not a refit.} The above
        additionally establishes that \texttt{ffsim}'s
        \texttt{interaction\_pairs} performs post-hoc zeroing of the
        factorised $J$ rather than a constrained fit to $t_2$. The analysis
        presupposes this, and it is now measured rather than assumed.
  \item \textbf{Localisation preserves the active space.} Localised and
        canonical CASCI energies agree to \num{8.5e-14}; the correlation
        energy reconstructed from rotated amplitudes against rotated
        integrals agrees with the canonical value to \num{2.3e-15}.
  \item \textbf{Bitstring convention.} Qiskit writes qubit 0 rightmost;
        \texttt{ffsim} maps alpha spin-orbitals to qubits $0..N-1$ and beta
        to $N..2N-1$; \texttt{sbd} expects the rightmost bit to be orbital 1.
        These coincide, so the required operation is a split, not a reversal.
\end{itemize}

\FloatBarrier
\section{Results}

\subsection{Orbital layout changes the answer by up to a factor of eight}

\begin{table}[htbp]\centering\small
\begin{tabular}{lrrrr}
\toprule
Ordering & Error (mHa) & Percentile & Captured & Retained $J$\\
\midrule
max.\ captured (oracle) & 24.27 & 0.7\% & 0.9831 & 0.2661\\
identity (inherited)    & 31.87 & 4.7\% & 0.9789 & 0.2874\\
max.\ retained $J$      & 44.32 & 22.1\%& 0.9668 & 0.3195\\
\midrule
\multicolumn{5}{l}{\emph{149 random layouts:} best 21.65, median 49.22, worst 173.23}\\
\bottomrule
\end{tabular}
\caption{\ce{N2}, CAS(6,10), 225-dimensional budget. Two seeds each;
seed-to-seed deviation zero throughout. A five-seed study gives between/within
variance ratio 27.2 and ranking stability $\tau=0.83$.}
\end{table}

\begin{table}[htbp]\centering\small
\begin{tabular}{lrr}
\toprule
Ordering & Error (mHa) & Percentile\\
\midrule
\texttt{s2\_max}            & 192.12 & 84\%\\
\texttt{physical\_reverse}  & 218.64 & 80\%\\
\texttt{s1\_max}            & 219.90 & 78\%\\
\texttt{retainedJ\_max}     & 295.82 & 26\%\\
reverse                     & 295.93 & 26\%\\
identity (inherited)        & 300.32 & 26\%\\
physical (nuclear sequence) & 389.71 & 2\%\\
\midrule
\multicolumn{3}{l}{\emph{50 random layouts:} best 168.67, median 263.19,
worst 454.89}\\
\bottomrule
\end{tabular}
\caption{\ce{H10} baseline. 114 evaluations, no support collapses,
seed-to-seed deviation zero for every named layout, between/within variance
ratio 278.2.}
\end{table}

\begin{figure}[htbp]\centering
\includegraphics[width=\linewidth,height=0.62\textheight,keepaspectratio]%
{fig1_effect_exists_paper.png}
\caption{The orbital-to-qubit assignment changes the SQD subspace error by up
to a factor of eight at fixed budget. Grey points are random layouts;
coloured markers are named layouts. Top: \ce{H10} CAS(10,10), $n=50$,
spanning \SIrange{168.67}{454.89}{\milli\hartree}. Bottom: \ce{N2}
CAS(6,10), $n=149$, spanning \SIrange{21.65}{173.23}{\milli\hartree}. Dotted
lines mark the random median. Seed-to-seed deviation is zero for every named
layout; between/within variance ratios are 278.2 and 27.2, so the spread is
signal rather than sampling noise.}
\label{fig:effect}
\end{figure}

The factor of eight is the full observed range on \ce{N2}. Between the fifth
and ninety-fifth percentiles the ratio is 2.71, and the largest-error layout
is not an outlier: it carries the lowest captured weight in the set, sitting
exactly where the mechanism of Section~\ref{sec:mechanism} predicts.
Excluding it moves the capture--error correlation from $-0.880$ to $-0.877$.

Two features of the \ce{H10} table matter more than the ranking. The
geometric ordering --- orbitals sequenced by localisation centroid along the
nuclear chain --- lands at the \emph{2nd} percentile. Spatial locality, the
most obvious chemically motivated criterion, is actively harmful here. And
\texttt{physical} differs from \texttt{physical\_reverse} by
\SI{171.07}{\milli\hartree}, 60\% of the observed random spread, despite
being the same ordering read backwards.

One \ce{N2} random layout produced fewer than fifteen unique determinants per
spin sector, collapsing the subspace to 81 dimensions with error near
\SI{206}{\milli\hartree}. Poor layouts can suppress sampling diversity below
the requested budget entirely.

\FloatBarrier
\subsection{The mechanism is subspace capture}
\label{sec:mechanism}

Defining captured weight as the fraction of the exact wavefunction inside the
sampled product subspace,
\begin{equation}
C = \!\!\sum_{a\in S_\alpha,\; b\in S_\beta}\!\! \lvert C_{ab}\rvert^{2},
\end{equation}
gives Spearman $\rho=-0.880$ ($p=\num{3e-49}$) across the 149 random \ce{N2}
layouts, $\rho=-0.983$ ($p=\num{3.6e-37}$) across the 50 random \ce{H10}
layouts, and $\rho=-0.97$ on \ce{Cr2}. The causal chain is
\begin{equation}
\text{layout} \rightarrow \text{sampled determinant distribution}
\rightarrow \text{captured weight at fixed budget}
\rightarrow \text{SQD subspace error}.
\end{equation}

\begin{figure}[htbp]\centering
\includegraphics[width=\linewidth,height=0.60\textheight,keepaspectratio]%
{fig2_mechanism_paper.png}
\caption{Subspace capture is the controlling variable, on all three systems.
Dashed vertical lines mark the ideal capture ceiling at each system's
determinant budget. \ce{N2} capture values cluster near the ceiling whereas
\ce{H10} spans \numrange{0.40}{0.71}, which is why approximate capture
surrogates fail on \ce{N2} for reasons of resolution rather than fidelity.}
\label{fig:mechanism}
\end{figure}

The relationship was subsequently tested on twelve same-spin chains held out
from all development (Section~\ref{sec:heldout}), where it holds at 12 of 12
without exception. Across the whole study it has been measured on twenty-one
distinct chains spanning three molecules, canonical and localised orbital
bases, weak and strong correlation regimes and both search axes, with $\rho$
between $-0.88$ and $-0.98$ throughout and no failures.

We emphasise that exact captured weight is an \emph{oracle diagnostic}. It
requires the exact wavefunction and is used to identify the mechanism at
classically tractable scale, not as a production selection criterion.

\FloatBarrier
\subsection{No scalar predictor of the permutation search works}

Eleven score variants --- amplitude-weighted retained-interaction measures and
dominant-excitation reachability measures constructed from CCSD $t_1,t_2$,
plus retained $J$ --- were evaluated against the 50-layout \ce{H10}
distribution. None reaches significance as a selector: $s_{1,\text{amp}}$
gives $\rho=-0.032$, $s_2$ gives $\rho=+0.091$, retained $J$ gives
$\rho=+0.001$. The strongest single variant reaches only $\rho=-0.377$
($p=\num{7.0e-3}$). Every selection regret equals or exceeds the
random-selection baseline of \SI{95.01}{\milli\hartree}.

Two structural explanations emerge, and both are more informative than the
null results themselves.

\textbf{Sign cancellation.} Retained $J$ decomposed by spin sector gives
$\rho=+0.212$ (not significant) for the same-spin contribution and
$\rho=-0.375$ ($p=\num{7.3e-3}$) for the opposite-spin contribution. The two
carry opposite signs and cancel to $+0.001$. The combined metric was
averaging a real signal against a spurious one --- and the surviving signal
lies entirely in the opposite-spin sector. This was the first indication that
the two sectors should be treated as separate degrees of freedom.

\textbf{Precision, not accuracy.} On \ce{N2}, capture values cluster in
\numrange{0.96}{0.98}, so differences below one per cent decide the ranking.
A CISD surrogate with marginal overlap 0.973 --- 97.3\% accurate --- reaches
only $\rho=-0.43$, because a target good to 2.7\% cannot resolve differences
of 0.5\%. Approximate capture targets fail on resolution, not fidelity.

\begin{figure}[htbp]\centering
\includegraphics[width=\linewidth,height=0.58\textheight,keepaspectratio]%
{fig6_null_results_paper.png}
\caption{No cheap scalar predicts the permutation search. Spearman $\rho$
against subspace error for all eleven score variants across the 50 random
\ce{H10} layouts, with bootstrap 95\% confidence intervals. Ten of eleven
intervals cross zero. The final two rows decompose retained $J$ by spin
sector, showing the cancellation.}
\label{fig:nulls}
\end{figure}

For completeness: mutual-information Fiedler ordering imported from DMRG
reaches the 27th percentile on \ce{N2}; an entanglement-weighted variant
reaches the 69th; and \emph{maximising} retained $J$ makes matters worse on
both systems. Orbital-ordering heuristics developed for
matrix-product-state methods, where ordering controls bond dimension, do not
transplant to this problem.

\FloatBarrier
\subsection{The mask has two separable levers}

The same-spin chain mask is invariant under reversal: orbitals adjacent at
positions $(k,k+1)$ remain adjacent at $(n-2-k,\,n-1-k)$. This was verified
symbolically for all 57 baseline permutations (57/57). Reversing a layout
therefore changes \emph{only} which orbitals occupy positions
$p\equiv 0\bmod 4$ and hence receive opposite-spin on-site coupling.

The \SI{171.07}{\milli\hartree} \texttt{physical}/\texttt{physical\_reverse}
gap is thus attributable entirely to anchor placement. Holding the same-spin
ordering fixed and sweeping the anchor offset over
$p\equiv\text{offset}\bmod 4$ spans \SI{171.08}{\milli\hartree}, agreeing
with the reversal gap to \SI{0.01}{\milli\hartree}. There is no universal
best offset: \texttt{identity} and \texttt{physical} prefer offset~1,
\texttt{physical\_reverse} and \texttt{s2\_max} prefer offset~0 --- an early
indication that the levers interact.

\begin{table}[htbp]\centering\small
\begin{tabular}{lrrr}
\toprule
Lever & Search space & Sampled & Observed range (mHa)\\
\midrule
Same-spin ordering (random permutations) & $10! = \num{3628800}$ & 50 & 286.23\\
Free anchor selection at fixed chain & $\binom{10}{3}=120$ & 120 & 234.10\\
Anchor offset at fixed chain & 4 & 4 & 171.08\\
\bottomrule
\end{tabular}
\caption{The anchor lever is comparable in magnitude to the permutation
search while spanning a space smaller by four orders of magnitude. Ranges are
observed over the stated samples at \ce{H10}; the anchor row is exhaustive
over all 120 triples at the identity chain, the permutation row a 50-element
sample of $10!$.}
\end{table}

\begin{figure}[htbp]\centering
\includegraphics[width=\linewidth,height=0.30\textheight,keepaspectratio]%
{fig3_two_levers_paper.png}
\caption{Observed error range attributable to each lever on \ce{H10},
annotated with the size of the corresponding search space.}
\label{fig:levers}
\end{figure}

\FloatBarrier
\subsection{Anchor optimisation compresses the layout spread}
\label{sec:compression}

To test whether the two levers interact, twenty same-spin chains spanning the
baseline range were each evaluated over a common set of anchor triples: eight
in an initial paired design, and twelve further chains held out from all
development.

\begin{table}[htbp]\centering\small
\begin{tabular}{lrrr}
\toprule
Chain & Default anchors & Best anchors & Change\\
\midrule
\texttt{rand030} (best baseline)  & 168.67 & 180.81 & $+12.14$\\
\texttt{rand037}                  & 175.60 & \textbf{170.47} & $-5.13$\\
\texttt{rand007}                  & 261.60 & 184.33 & $-77.27$\\
\texttt{rand032}                  & 264.78 & 172.97 & $-91.81$\\
\texttt{identity} (inherited)     & 300.32 & 230.31 & $-70.01$\\
\texttt{rand047}                  & 383.16 & 181.38 & $-201.78$\\
\texttt{physical}                 & 389.71 & 175.17 & $-214.54$\\
\texttt{rand029} (worst baseline) & 454.89 & 186.74 & $-268.15$\\
\bottomrule
\end{tabular}
\caption{The initial eight-chain paired design. \texttt{rand030}'s
best-of-candidates exceeds its default because its default anchor triple was
not among the sampled candidates; a comparability audit found this to be the
case for 6 of these 8 chains, so this table understates the achievable
improvement. The compression factor quoted below is computed over twenty
chains with the default triple always included.}
\end{table}

Three features matter. First, the \textbf{compression}: most of what presents
as an orbital-ordering effect is an anchor-placement effect. The same-spin
permutation and the opposite-spin anchor set are not independent
contributions but largely substitutable ones, and a poor permutation can be
substantially compensated by a good anchor choice.

Second, the \textbf{rank reversals}. \texttt{rand029}, the worst of all 50
random layouts at default anchors, attains \SI{186.74}{\milli\hartree} once
its anchors are optimised --- better than \texttt{rand030}, the best of all
50, attains at \SI{180.81}{\milli\hartree}. \texttt{physical}, at the 2nd
percentile of the baseline distribution, reaches
\SI{175.17}{\milli\hartree}. The global optimum in the eight-chain design
came from a chain at only the 75th percentile.

Third, \textbf{residual chain dependence}. The post-optimisation spread is
not zero. Same-spin ordering continues to matter; it matters far less than
the raw baseline distribution suggests. The rank correlation between baseline
and post-optimisation error was $\rho = +0.405$ ($p = 0.320$) on the initial
eight chains --- inconclusive rather than null at that sample size. We
therefore do not claim that baseline ranking fails to predict
post-optimisation ranking, only that the compression and the observed
reversals are inconsistent with baseline quality being a reliable guide.

\begin{figure}[htbp]\centering
\includegraphics[width=\linewidth,height=0.45\textheight,keepaspectratio]%
{fig7a_compression_paper.png}
\caption{Anchor optimisation compresses the layout spread. Each line connects
one same-spin chain's subspace error at the inherited $p\equiv 0\bmod 4$
anchor placement to its error at the best anchor triple found. The fan
collapses and the lines cross: the worst baseline chain overtakes the best.}
\label{fig:interaction}
\end{figure}

The practical implication is direct. Optimising the same-spin permutation
first and accepting whatever anchor placement follows --- the conventional
framing, and our own until this experiment --- addresses the smaller and far
more expensive lever.

\FloatBarrier
\subsection{The anchor effect is present in the ansatz itself}
\label{sec:ansatz}

Every result above measures SQD subspace error, which depends jointly on the
ansatz, the sampling and the determinant selection. To separate these we
computed the variational energy of the masked LUCJ state directly,
\begin{equation}
\varepsilon_{\text{LUCJ}} = 10^{3}
\left(
\frac{\langle\psi_{\text{LUCJ}}\lvert \hat H \rvert\psi_{\text{LUCJ}}\rangle}
{\langle\psi_{\text{LUCJ}}\vert\psi_{\text{LUCJ}}\rangle}
- E_{\text{CASCI}}
\right),
\end{equation}
in the number-conserving sector, with no sampling and no diagonalisation.

The anchor effect is clearly present at the ansatz level. At \ce{H10}
identity, $\varepsilon_{\text{LUCJ}}$ spans \SI{135.03}{\milli\hartree}
against the \SI{234.10}{\milli\hartree} SQD range, a ratio of 0.577; the ratio
stays in \numrange{0.44}{0.63} at every chain. The
``anchors can be worse than nothing'' result of
Section~\ref{sec:floor} also replicates independently: 5 of 120 triples
(4.2\%) give a worse variational energy than the no-$\alpha\beta$ control,
against 4.0\% measured at the SQD level in an entirely separate pipeline.

Ansatz quality transmits to the SQD outcome only partially, and only at some
chains. At \ce{H10} identity,
$\rho(\varepsilon_{\text{LUCJ}}, \varepsilon_{\text{SQD}}) = +0.527$, and
$+0.475$ ($p=\num{4.1e-7}$) with the 17 degenerate floor triples excluded, so
the relationship is not an artefact of the extremes. At \texttt{physical} and
\texttt{rand007} it is not significant either way. This is a direct
demonstration of the VQE/SQD distinction: a better ansatz is not reliably a
better sampler.

\FloatBarrier
\subsection{The no-coupling floor, and anchors that actively harm}
\label{sec:floor}

Seventeen of the 120 triples at \ce{H10} identity return an identical error
of \SI{458.70}{\milli\hartree}. A direct control with the opposite-spin
interaction list emptied returns \SI{458.6997}{\milli\hartree}, identifying
this as the \emph{no-$\alpha\beta$-correlation floor}: those triples place
anchors on orbitals whose $\lvert J^{\alpha\beta}_{pp}\rvert$ is negligible
(mean 1.30, against 8.03 for the best ten), so the opposite-spin sector
contributes nothing.

\begin{table}[htbp]\centering\small
\begin{tabular}{lrrrr}
\toprule
Chain & Floor (no $\alpha\beta$) & Default anchors & Best anchors & Worse than floor\\
\midrule
identity  & 458.70 & 300.32 & 224.60 & 0/120\\
rand007   & 372.86 & 261.60 & 175.54 & 3/40\\
physical  & 379.25 & \textbf{389.71} & \textbf{172.15} & 1/40\\
\bottomrule
\end{tabular}
\caption{At \texttt{physical} the inherited $p\equiv 0\bmod 4$ anchor
placement is \SI{10.46}{\milli\hartree} \emph{worse} than removing
opposite-spin coupling altogether. Both figures are 15-dimensional at
$2\times10^{6}$ shots and reproduce exactly on re-run.}
\end{table}

\begin{figure}[htbp]\centering
\includegraphics[width=\linewidth,height=0.45\textheight,keepaspectratio]%
{fig5_floor_and_limits_paper.png}
\caption{The no-$\alpha\beta$-correlation floor, the inherited anchor
placement and the best anchor selection found, at three same-spin chains.
Dashed lines carry each chain's own floor across its group.}
\label{fig:floor}
\end{figure}

Across all 50 random layouts, 2 of 50 (4.0\%) have a default-anchor error
worse than their own floor, with a mean penalty of
\SI{33.75}{\milli\hartree}. The effect is a minority case but systematic:
$\rho(\text{baseline error},\; \text{floor}-\text{default gap}) = -0.359$
($p=0.010$), meaning opposite-spin coupling helps at good chains and hurts at
bad ones. A guard that disables the opposite-spin sector when it fails to
help costs one additional evaluation per chain.

An apparently larger anomaly on \ce{N2} --- 20\% of triples worse than the
floor at one chain --- proved to be floor proximity rather than structure.
That chain's floor lies at \SI{49.69}{\milli\hartree} against
\SI{114.00}{\milli\hartree} elsewhere, so an equal absolute penalty crosses it
far more often; the worse-than-floor triples show no distinguishing signature
in any diagnostic (all Mann--Whitney $p>0.27$).

\FloatBarrier
\subsection{Transfer across systems}
\label{sec:transfer}

The anchor axis was repeated on \ce{N2} --- a different molecule, canonical
rather than localised orbitals, weak correlation --- and on \ce{Cr2}, a 3d
transition-metal dimer with a localised active space.

\begin{table}[htbp]\centering\small
\begin{tabular}{lrrr}
\toprule
& \ce{H10} & \ce{N2} & \ce{Cr2}\\
\midrule
Anchor range at identity (mHa) & 234.10 & 89.73 & ---\\
Capture ceiling & 0.7554 & 0.9866 & 0.881\\
Headroom (ceiling $-$ mean achieved) & 0.1620 & 0.0352 & ---\\
\textbf{Range per unit headroom (mHa)} & \textbf{1445} & \textbf{2550} & \textbf{1677}\\
$\rho(C, \varepsilon_{\text{SQD}})$ & $-0.983$ & $-0.877$ & $-0.97$\\
$\rho(S_0, \varepsilon_{\text{LUCJ}})$ & $-0.850$ & $-0.965$ & $-0.96$ to $-0.97$\\
Best triple moves between chains & yes & yes & yes\\
\bottomrule
\end{tabular}
\caption{Anchor-axis quantities on all three systems. The raw anchor range is
smaller on \ce{N2}, but normalised by available capture headroom the \ce{N2}
lever is the \emph{strongest} of the three: the smaller raw number reflects
proximity to the ceiling, not a weaker effect.}
\end{table}

Two observations. The headroom normalisation, introduced to explain the
\ce{N2} result, is supported by a third independent point on \ce{Cr2}, which
falls between the other two. And at \ce{Cr2}, identity and reverse chains
gave bit-identical results, independently confirming the reversal-invariance
argument on a system from which it was never derived.

\FloatBarrier
\subsection{Why the chain dependence exists}
\label{sec:transmission}

Decomposing the pathway into two links,
\begin{equation}
\underbrace{\varepsilon_{\text{LUCJ}} \rightarrow C}_{\text{link 1}}
\qquad
\underbrace{C \rightarrow \varepsilon_{\text{SQD}}}_{\text{link 2}},
\end{equation}
and measuring each at six chains across two systems, link 2 holds everywhere
without exception ($\rho$ between $-0.88$ and $-0.96$). Link 1 fails at
\ce{H10} \texttt{physical} and \texttt{rand007}. The bottleneck is not the
subspace diagonalisation but whether a variationally good state also
concentrates its sampling distribution usefully. Six candidate diagnostics of
the sampled distribution --- unique string counts, entropy, Gini coefficient,
top-15 mass, and the boundary ratio --- fail to explain the residual
consistently; the best is significant at only 2 of 6 chains.

\paragraph{No purely opposite-spin rule can select the optimum at every chain.}
Define
\begin{equation}
S_0(A) \;=\; \sum_{p\in A} \bigl\lvert J^{\alpha\beta}_{pp}\bigr\rvert,
\label{eq:s0}
\end{equation}
summed over repetitions and normalised, where $A$ is the anchor set. The
decisive observation is algebraic rather than statistical: $S_0$ depends only
on \emph{which} orbitals are anchored, never on the same-spin permutation. It
is therefore exactly invariant under that permutation and assigns an
identical set of values to the candidate anchor sets at every chain. The
chain-specific optimum, however, moves: the best triple differs at every
chain tested on all three systems. A stationary ranking is being asked to
track a moving target.

This sets a precise requirement on any successor: \emph{a rule that selects
the optimal anchor set at every chain must incorporate same-spin chain
information}. Purely opposite-spin quantities are provably insufficient for
that objective, however well they perform at a particular chain.

\FloatBarrier
\subsection{Three families of chain-aware alternatives, all pre-registered, all failed}
\label{sec:chainaware}

Acting on that requirement, three families of chain-aware scores were
declared in advance of evaluation across three separate attempts, with
development restricted to identity chains and all other chains held out.

\textbf{Round 1 --- four scores targeting the variational energy.}
$S_1$ reach-weighted, multiplying each anchor's opposite-spin weight by the
retained same-spin coupling strength in its chain neighbourhood; $S_2$
coverage, penalising anchors that cluster on the chain; $S_3$
amplitude-reach; $S_4$ combining $S_1$ and $S_2$. One declared hyperparameter
sweep. None improved on $S_0$ in worst-case correlation.

\textbf{Round 2 --- two scores targeting capture directly.}
$T_1$ anchor-conditioned reachability of dominant CCSD excitation channels;
$T_2$ a perturbative estimate of the masked state's support. Both were
markedly \emph{worse} than $S_0$: their worst-case correlations with capture
are $-0.486$ and $-0.644$ against $S_0$'s $+0.361$, and they carry the wrong
sign at 14 and 15 of 18 chains respectively. Targeting capture directly
performs worse than targeting the ansatz energy.

\textbf{Round 3 --- an anchor--neighbour interface score.}
\begin{equation}
S_{\text{int}}(c, A) = \sum_{a \in A}\;
\sum_{n \in N_{\text{same}}(a;c)}
\bigl\lvert J^{\alpha\beta}_{an}\bigr\rvert \cdot
\bigl\lvert J^{\text{same}}_{an}\bigr\rvert,
\end{equation}
with zero free parameters, on the rationale that an anchor seeds
$\alpha\beta$ correlation which the chain then propagates through retained
same-spin couplings. Unlike rounds 1 and 2 this passes an automatic
invariance pre-flight --- its value genuinely changes under same-spin
permutation --- but it does not improve on $S_0$ under the shortlist
criterion of Section~\ref{sec:shortlist} (median regret 0.036 against
$S_0$'s 0.002).

Seven pre-declared scores across three attempts and two prediction targets,
none improving on the simple opposite-spin quantity. We report this as a
closed negative result: the obvious ways of injecting same-spin information
into an anchor score do not repair the chain dependence.

\FloatBarrier
\subsection{Out-of-sample validation on twelve held-out chains}
\label{sec:heldout}

Twelve \ce{H10} same-spin chains were drawn from a fixed seed and verified
disjoint from every chain used anywhere in the study. All 120 anchor triples
were evaluated at the ansatz level at each (1440 evaluations), and a paired
43-triple set plus the default placement and the no-$\alpha\beta$ control at
the SQD level (540 evaluations, zero support collapses).

\begin{table}[htbp]\centering\small
\begin{tabular}{lr}
\toprule
Quantity across 18 chains & $S_0$\\
\midrule
$\rho(S_0, C)$: correct sign & \textbf{18/18}\\
$\rho(S_0, C)$: significant at $p<0.05$ & \textbf{18/18}\\
$\rho(S_0, C)$: median & $+0.703$\\
$\rho(S_0, C)$: worst case & $+0.361$\\
$\rho(S_0, \varepsilon_{\text{SQD}})$: correct sign & \textbf{18/18}\\
Argmax normalised regret: median & 0.347\\
Argmax normalised regret: chains exceeding 1.0 & 1/18\\
$\rho(C, \varepsilon_{\text{SQD}})$ holds & \textbf{12/12} held out\\
\bottomrule
\end{tabular}
\caption{$S_0$ correlates with capture in the correct direction, significantly,
at every one of eighteen chains. Its single argmax failure
(\texttt{newchain11}, regret 1.152) traces to an exact tie in the score
between two candidate triples, broken by an uninformative lexicographic rule;
the score's correlation at that chain is unaffected ($+0.609$, significant).}
\end{table}

\FloatBarrier
\subsection{The deployable objective is a shortlist, not a single pick}
\label{sec:shortlist}

Selecting one anchor set blind by $\mathrm{argmax}\,S_0$ is not how the rule
would be used. With 120 candidates and an affordable evaluation, a
practitioner ranks them and evaluates a handful. We therefore report
\emph{best-of-top-$k$} normalised regret: rank all candidates by the score,
evaluate the top $k$, and take the best result.

\begin{table}[htbp]\centering\small
\begin{tabular}{lrrrr}
\toprule
Best-of-top-5 normalised regret & $S_0$ & $S_{\text{int}}$ & $S_0+S_{\text{int}}$ & random\\
\midrule
median & \textbf{0.002} & 0.036 & 0.008 & 0.181\\
90th percentile & \textbf{0.002} & 0.036 & 0.008 & 0.621\\
\bottomrule
\end{tabular}
\caption{Twelve held-out chains, evaluated once. Ranking the 120 anchor sets
by $S_0$ and evaluating only the top five reaches within 0.2\% of the best
candidate's error, with no tail: the 90th percentile equals the median.
Random selection of five candidates gives 0.181 and 0.621.}
\end{table}

This is the practical resolution of the tension in
Section~\ref{sec:transmission}. $S_0$ provably cannot identify the optimum at
every chain, and its argmax does fail at one chain in eighteen. But under
shortlisting the invariance limitation costs \SI{0.2}{\percent} of the
available improvement, at a cost of five evaluations from a space of 120.

\FloatBarrier
\subsection{Circuit resource accounting}
\label{sec:resources}

All comparisons above hold fixed the mask size, shot budget and determinant
budget. Whether they also hold circuit cost fixed was verified directly, by
transpiling each layout's masked LUCJ circuit against a heavy-hex coupling
map (\texttt{CouplingMap.from\_heavy\_hex}, distance 5, 57 qubits) with five
independent routing seeds per layout.

\begin{table}[htbp]\centering\small
\begin{tabular}{lrrrr}
\toprule
Configuration & $\Delta$ error (mHa) & $\Delta$ 2Q gates & gates / mHa\\
\midrule
\ce{H10} identity, default $\to$ best & $-75.72$ & $-66.6$ & $-0.88$\\
\ce{H10} physical, default $\to$ best & $-217.56$ & $+336.6$ & $+1.55$\\
\ce{N2} identity, default $\to$ best & $-7.60$ & $-880.8$ & $-115.89$\\
\ce{Cr2} identity, default $\to$ best & $-38.84$ & $-2823.4$ & $-72.69$\\
\bottomrule
\end{tabular}
\caption{Resource cost of each quoted improvement. Three of the four recover
accuracy while \emph{reducing} two-qubit gate count. Only \ce{H10}
\texttt{physical} pays a resource cost, and it is small relative to that
circuit's total size.}
\end{table}

Layouts are \emph{not} exactly resource-neutral. Across the 120 anchor
triples at \ce{H10} identity the coefficient of variation is 5.1\% in
two-qubit gate count, 11.6\% in depth and 27.4\% in SWAP count; across 20
same-spin chains it is 19.0\%, 21.7\% and 37.5\%. Layout variation exceeds
routing-seed variation by a factor of 1.39, so it is signal rather than
noise. The 5.1\% figure on the anchor axis is 96\% attributable to routing:
the logical, pre-routing two-qubit gate count varies by only 0.27\% across
the 120 triples, as expected given that every triple retains exactly three
opposite-spin terms.

Two findings favour the anchor axis. On that axis, resource cost and accuracy
move \emph{together}: $\rho(\text{2Q gates}, S_0) = -0.709$ and
$\rho(\text{2Q gates}, \varepsilon_{\text{SQD}}) = +0.608$, so better anchor
sets require fewer gates. And the same-spin axis is by far the dominant
hardware cost: \texttt{physical} requires roughly 26{,}500 two-qubit gates
and 5{,}284 SWAPs against \texttt{identity}'s 12{,}800 and 650, a factor of
two in gates and eight in SWAPs.

We therefore do not claim ``fixed quantum resources'' in the strict sense.
The defensible statement is that anchor selection improves accuracy at
approximately constant, and typically slightly reduced, two-qubit gate count,
while same-spin reordering can double it.

\FloatBarrier
\section{Deployment and scaling}

The scaling argument is the practical case for the anchor axis. At
$N_{\text{orb}}=20$ with five anchor positions the anchor space is
$\binom{20}{5}=\num{15504}$; at $N_{\text{orb}}=36$ with nine it is
$\binom{36}{9}\approx\num{9.4e7}$, against $36!\approx\num{3.7e41}$ for the
permutation search. Combined with the shortlist result of
Section~\ref{sec:shortlist}, the deployment recipe is: rank the anchor
candidates by $S_0$, evaluate the top five, take the best. On the twelve
held-out chains that recipe lands within 0.2\% of the best available
candidate.

An A64FX build of \texttt{sbd} has been established on Fugaku via the
repository's LLVM/OpenBLAS CMake preset, cross-compiled with
\texttt{mpiclang++} and verified as an \texttt{aarch64} binary, with job
submission through \texttt{pjsub} confirmed working end to end. Numerical
validation against the group benchmark is incomplete and no Fugaku-derived
numbers appear in this report. Timing measurements indicate that Fugaku's
value for this workload is throughput rather than single-job latency: the
anchor sweep is embarrassingly parallel across many short independent
diagonalisations, whereas individual A64FX cores are slower than the local
development hardware.

The binding constraint on system size is classical simulation, not the
solver. Each \ce{Cr2} evaluation at 24 qubits required roughly 16 minutes, of
which 962 seconds was Aer statevector simulation; cost scales as
$2^{2N_{\text{orb}}}$. A [2Fe--2S]-scale active space at 40 qubits is
therefore out of reach of this workflow classically, though on hardware the
sampling would be free and the entire cost would be the anchor enumeration
--- precisely the axis shown here to be small.

\FloatBarrier
\section{Contributions}

\begin{enumerate}\itemsep3pt
  \item A validated pipeline reproducing the group's published benchmark to
        twelve significant figures, with end-to-end permutation-invariance
        and cross-implementation equivalence checks.
  \item Quantification of the orbital-layout effect on two systems, robust to
        sampling noise (between/within variance ratios 27.2 and 278.2).
  \item Identification of subspace capture as the governing mechanism,
        holding at twenty-one chains across three molecules, two orbital
        bases, two correlation regimes and both search axes, including twelve
        chains held out from all development.
  \item A characterised negative result on cheap prediction of the
        permutation search --- eleven variants --- with two structural
        explanations: spin-sector sign cancellation, and a resolution rather
        than fidelity limit on approximate capture targets.
  \item \textbf{Decomposition of the locality mask into two separable but
        interacting levers}, with the opposite-spin anchor lever shown
        comparable in magnitude to the permutation search while occupying a
        space four orders of magnitude smaller.
  \item Demonstration that anchor optimisation compresses the observed layout
        spread across twenty chains and reverses their ranking.
  \item \textbf{A proof that no purely opposite-spin quantity can select the
        optimal anchor set at every chain}, with seven pre-registered
        chain-aware alternatives across three attempts that fail to repair
        it.
  \item \textbf{A deployable shortlist rule}: ranking anchor candidates by
        opposite-spin retained Jastrow weight and evaluating the top five
        reaches median and 90th-percentile normalised regret of 0.002 on
        twelve held-out chains, against 0.181 and 0.621 for random selection.
  \item A circuit-resource audit establishing that better anchor sets require
        fewer two-qubit gates, and that same-spin reordering rather than
        anchor selection is the dominant hardware cost.
\end{enumerate}

\section{Limitations and non-claims}

\begin{itemize}\itemsep3pt
  \item We do \emph{not} claim a rule that identifies the optimal anchor set
        at every chain. The opposite is proved for opposite-spin-only
        quantities.
  \item We do \emph{not} claim strict fixed quantum resources. Layouts differ
        in transpiled gate count and depth (Section~\ref{sec:resources});
        comparisons are at fixed mask size, shot budget and determinant
        budget.
  \item Captured weight is an oracle diagnostic requiring the exact
        wavefunction. It identifies the mechanism; it is not a production
        criterion.
  \item Orbital layout does not repair the failure of a single-reference
        ansatz on a strongly correlated system. The \ce{Cr2} study is a
        transfer test of the layout phenomenon, not a chemistry result.
  \item Observed error ranges are over the stated samples, not bounds over
        the full search space.
  \item All results are at a single shot budget per system and a single
        determinant budget; no hardware execution was performed.
\end{itemize}

\section{Next steps}

\begin{enumerate}\itemsep3pt
  \item \textbf{A chain-aware rule with a generality guarantee}, beyond the
        three families tested. The requirement is precisely specified by
        Section~\ref{sec:transmission}; what remains open is whether any
        cheap quantity satisfies it.
  \item \textbf{Explain link-1 failure directly}, comparing the
        $J^{\alpha\beta}_{pp}$ distribution, occupied/virtual character and
        per-anchor capture contributions between chains where ansatz quality
        transmits and where it does not.
  \item \textbf{Validate the shortlist recipe on a fourth system} and at a
        second shot budget, to establish that the 0.002 regret figure is a
        property of the method rather than of \ce{H10}.
  \item \textbf{Complete Fugaku numerical validation} and demonstrate the
        job-array route, so anchor enumeration at $N_{\text{orb}}\ge 20$
        becomes routine.
  \item \textbf{Hardware execution}, where sampling is free and the anchor
        enumeration is the entire cost.
\end{enumerate}

\FloatBarrier
\section*{Reproducibility}

All code, cached reference data, per-evaluation result files and experiment
metadata are available at
\url{https://github.com/solovevmax/sqd-orbital-ordering}. Each experiment
directory records the git commit, SHA-256 hashes of every reference artefact
used, shot counts, random seeds and wall time. The \texttt{sbd} solver is used
unmodified throughout.

\paragraph{A withdrawn result, and why a rank statistic missed it.}
An earlier reimplementation of the locality mask omitted the same-spin
diagonal $(p,p)$ terms --- 10 of 31 retained entries per repetition. All
\ce{H10} results obtained before 25 August 2026 were consequently withdrawn
and re-measured. \ce{N2} results were unaffected: the original implementation
always included the diagonal.

The omission is physically consequential. Since
$J_{pp}n_{p\sigma}n_{p\sigma} = J_{pp}n_{p\sigma}$ by idempotency, these
terms are single-qubit $Z$ rotations, unconstrained by heavy-hex connectivity
and therefore never a legitimate target for a locality mask.

The validation that failed to detect it is instructive. The two mask models
were compared by Spearman correlation of retained $J$ and agreed at $+0.993$.
But the diagonal is permutation-invariant: omitting it subtracts the same
constant from every layout, and a rank statistic is by construction blind to
an additive offset. A Pearson comparison on absolute values, or a comparison
of retained-pair counts, would have caught it immediately. Both
implementations now import a single shared mask module, guarded by a
regression test over 20 random permutations.

\paragraph{A corrected compression figure.} An earlier eight-chain estimate of
the anchor-optimisation compression factor was superseded when a
comparability audit established that the default anchor triple was absent
from the sampled candidate set at 6 of those 8 chains, so ``best of
candidates'' had excluded the incumbent. The figure quoted in
Section~\ref{sec:compression} is recomputed over twenty chains with the
default always included.

\paragraph{A corrected verdict.} A sign-convention error in the score-comparison
analysis silently substituted each score's \emph{best} chain for its
\emph{worst}, producing two false ``beats $S_0$'' verdicts in round 2. The
per-chain tables, hyperparameter sweep and regret figures were correct
throughout; only the worst-case summary statistic and the verdict derived
from it were affected. The corrected result is reported in
Section~\ref{sec:chainaware}.

\end{document}
